\documentclass[12pt,a4paper]{article}

\usepackage[a4paper,margin=1in,headheight=15pt]{geometry}
\usepackage{graphicx}
\usepackage{booktabs}
\usepackage{amsmath,amssymb}
\usepackage{float}
\usepackage{hyperref}
\usepackage{xcolor}
\usepackage{listings}
\usepackage{enumitem}
\usepackage{fancyhdr}
\usepackage{longtable}
\usepackage{array}
\usepackage{multirow}
\usepackage{subcaption}
\usepackage{caption}
\usepackage{url}
\usepackage{xurl}
\usepackage{hyperref}

\Urlmuskip=0mu plus 2mu\relax
\sloppy

\hypersetup{colorlinks=true,linkcolor=blue,urlcolor=blue}

\pagestyle{fancy}
\fancyhf{}
\lhead{ICS1512}
\rhead{Machine Learning Algorithms Laboratory}
\cfoot{\thepage}

\definecolor{codebg}{RGB}{245,245,245}
\definecolor{keywordcolor}{RGB}{0,0,180}
\definecolor{commentcolor}{RGB}{80,120,80}
\definecolor{stringcolor}{RGB}{163,21,21}

\lstset{
  language=Python,
  basicstyle=\ttfamily\footnotesize,
  numbers=left,
  numberstyle=\tiny\color{gray},
  frame=single,
  breaklines=true,
  showstringspaces=false,
  backgroundcolor=\color{codebg},
  keywordstyle=\color{keywordcolor}\bfseries,
  commentstyle=\color{commentcolor}\itshape,
  stringstyle=\color{stringcolor},
  tabsize=4,
  captionpos=b
}

\begin{document}

% ─────────────────────────────────────────────────────────────────────────────
%  COVER TABLE
% ─────────────────────────────────────────────────────────────────────────────
\begin{center}
  \textbf{\large Sri Sivasubramaniya Nadar College of Engineering, Chennai}\\[4pt]
  \textit{An Autonomous Institution Affiliated to Anna University}\\[12pt]
\end{center}

\begin{tabular}{|p{4.5cm}|p{10.5cm}|}
\hline
\textbf{Experiment No.}   & 5 \\ \hline
\textbf{Experiment Title} & Decision Tree and Random Forest: A Comparative Classification Study \\ \hline
\textbf{GitHub Repository} & \url{https://github.com/RadeshL/Machine-Learning-Laboratory/blob/main/MLassign5.ipynb} \\ \hline
\textbf{Degree \& Branch} & M.Tech (Integrated) Computer Science \& Engineering \\ \hline
\textbf{Semester}         & V \\ \hline
\textbf{Subject Code \& Name} & ICS1512 -- Machine Learning Algorithms Laboratory \\ \hline
\textbf{Academic Year}    & 2026--2027 (Odd Semester) \\ \hline
\textbf{Batch}            & 2024--2029 \\ \hline
\textbf{Student Name}     & L.Radesh\\ \hline
\textbf{Register Number}  & 3122247001047\\ \hline
\textbf{Faculty}          & Dr. Poreddy Ajay Kumar Reddy \\ \hline
\textbf{Submission Date}  & 16-08-2026 \\ \hline
\end{tabular}

\vspace{1cm}

% ─────────────────────────────────────────────────────────────────────────────
\section{Objective}
% ─────────────────────────────────────────────────────────────────────────────
\begin{itemize}[noitemsep]
  \item To implement a Decision Tree classifier on the Wisconsin Diagnostic Breast Cancer dataset.
  \item To extend the single Decision Tree into a Random Forest ensemble model.
  \item To study the impact of hyperparameters on overfitting and generalisation.
  \item To select optimal hyperparameters using 5-Fold Stratified Cross-Validation.
  \item To compare single-tree and ensemble-tree models using standard evaluation metrics.
\end{itemize}

% ─────────────────────────────────────────────────────────────────────────────
\section{Problem Statement}
% ─────────────────────────────────────────────────────────────────────────────
Given a set of 30 real-valued features derived from digitised fine-needle aspirate (FNA) images of breast
masses, the task is to classify each sample as \textbf{Malignant} or \textbf{Benign}.
The input is a feature vector $\mathbf{x} \in \mathbb{R}^{30}$ and the output is a binary label
$y \in \{0 (\text{Malignant}),\; 1 (\text{Benign})\}$.
Two tree-based classifiers---a single Decision Tree and a Random Forest ensemble---are trained,
tuned via 5-fold cross-validation, and compared to determine which achieves better diagnostic accuracy
while maintaining clinical-grade recall for the malignant class.

% ─────────────────────────────────────────────────────────────────────────────
\section{Dataset Description}
% ─────────────────────────────────────────────────────────────────────────────

\begin{longtable}{|p{5.5cm}|p{9cm}|}
\hline
\textbf{Attribute}      & \textbf{Detail} \\ \hline
Dataset Name            & Wisconsin Diagnostic Breast Cancer (WDBC) \\ \hline
Dataset Source          & UCI Machine Learning Repository (\url{https://archive.ics.uci.edu}) \\ \hline
Number of Samples       & 569 \\ \hline
Number of Features      & 30 numerical attributes (mean, SE, worst of 10 cell-nucleus measurements) \\ \hline
Number of Classes       & 2 --- Malignant (M): 212, Benign (B): 357 \\ \hline
Missing Values          & None \\ \hline
Train--Test Split       & 80\% training (455 samples) / 20\% test (114 samples), stratified \\ \hline
Label Encoding          & Malignant $\rightarrow$ 0, Benign $\rightarrow$ 1 \\ \hline
\end{longtable}

% ─────────────────────────────────────────────────────────────────────────────
\section{Exploratory Data Analysis}
% ─────────────────────────────────────────────────────────────────────────────

\subsection*{4.1 Class Distribution}

\begin{figure}[H]
\centering
\includegraphics[width=0.52\linewidth]{fig1_class_dist}
\caption{Class distribution of the WDBC dataset.
  The dataset contains 212 malignant and 357 benign samples, yielding a
  class imbalance ratio of approximately 1:1.68.}
\end{figure}

\textbf{Inference:}
The dataset exhibits a moderate class imbalance with benign samples outnumbering malignant ones by a
ratio of roughly 1.68:1. This imbalance is not severe enough to require resampling techniques, but it
motivates the use of stratified splitting and F1-score alongside accuracy as evaluation metrics.
The minority class (Malignant) is the clinically critical one---missed malignant predictions carry a
higher cost than false positives---making recall for class 0 a key diagnostic concern.

\subsection*{4.2 Feature Correlation Analysis}

\begin{figure}[H]
\centering
\includegraphics[width=0.88\linewidth]{fig2_correlation}
\caption{Pairwise Pearson correlation heatmap for the 12 highest-variance features.
  Dark red indicates strong positive correlation; dark blue indicates strong negative correlation.}
\end{figure}

\textbf{Inference:}
Several feature groups exhibit very high mutual correlation (|r| $>$ 0.9), particularly the
\textit{radius}, \textit{perimeter}, and \textit{area} features across the mean, SE, and worst
measurement groups---which is geometrically expected since these quantities are algebraically related
($A = \pi r^2$). This multicollinearity has limited impact on tree-based methods (which select splits
greedily), but it does make feature importance scores somewhat redundant across correlated groups.
The strongest diagnostic signal is concentrated in the \textit{worst} and \textit{mean} concavity
and concave-points features, consistent with the literature on FNA morphology.

% ─────────────────────────────────────────────────────────────────────────────
\section{Data Preprocessing}
% ─────────────────────────────────────────────────────────────────────────────

\begin{itemize}[noitemsep]
  \item \textbf{Missing values:} None present; no imputation required.
  \item \textbf{Label encoding:} String labels M/B from the original UCI file are mapped to
        integers (M$\rightarrow$0, B$\rightarrow$1) by \texttt{sklearn}'s built-in dataset loader.
  \item \textbf{Feature scaling:} Not applied. Decision Trees and Random Forests are
        scale-invariant; their split thresholds operate on raw feature magnitudes.
  \item \textbf{Train--test split:} 80--20 stratified split with \texttt{random\_state=42},
        preserving the class ratio in both partitions.
  \item \textbf{Feature engineering:} No derived features added; all 30 raw numerical attributes
        are used directly.
\end{itemize}

% ─────────────────────────────────────────────────────────────────────────────
\section{Machine Learning Algorithm}
% ─────────────────────────────────────────────────────────────────────────────

\subsection*{6.1 Decision Tree Classifier}

\textbf{Working Principle:}
A Decision Tree recursively partitions the feature space by selecting at each node the feature and
threshold that maximally reduce impurity in the resulting child nodes, continuing until a stopping
criterion (maximum depth, minimum samples) is met. Prediction for a new sample traces a path from
root to leaf following the learned rules.

\textbf{Mathematical Formulation:}
At each node $t$, the best split is chosen by minimising the weighted child impurity:

\[
\Delta I(t) = I(t) - \frac{N_L}{N} I(t_L) - \frac{N_R}{N} I(t_R)
\]

where $N_L, N_R$ are child node sample counts and $I(\cdot)$ is the impurity measure. For
\textbf{Gini index}:
\[
I_{\text{Gini}}(t) = 1 - \sum_{k=1}^{K} p_k^2
\]
and for \textbf{Entropy}:
\[
I_{\text{Entropy}}(t) = -\sum_{k=1}^{K} p_k \log_2 p_k
\]
where $p_k$ is the proportion of class $k$ samples at node $t$.

\textbf{Advantages:} Highly interpretable; no feature scaling needed; handles mixed feature types.

\textbf{Limitations:} High variance; prone to overfitting on deep trees; sensitive to small changes
in training data.

\subsection*{6.2 Random Forest Classifier}

\textbf{Working Principle:}
Random Forest constructs $T$ decision trees, each trained on a bootstrapped resample of the
training data. At each split, only a random subset of $m$ features is considered.
Final prediction is determined by majority vote across all trees.

\textbf{Mathematical Formulation:}
For $T$ trees $\{h_1, h_2, \ldots, h_T\}$, the ensemble prediction is:
\[
\hat{y} = \underset{k}{\arg\max} \sum_{t=1}^{T} \mathbf{1}[h_t(\mathbf{x}) = k]
\]
Bootstrap samples are drawn with replacement: each subsample $\mathcal{D}_t \sim \mathcal{D}$
with $|\mathcal{D}_t| = N$. Feature subspace size $m = \lfloor \log_2 p \rfloor$ (\texttt{log2})
for classification, where $p$ is the total feature count.

The expected generalisation error satisfies:
\[
\text{PE}^* \leq \frac{\bar{\rho}(1 - s^2)}{s^2}
\]
where $\bar{\rho}$ is the mean inter-tree correlation and $s$ is the mean tree strength (Breiman, 2001).
Reducing $\bar{\rho}$ via random feature selection is the primary mechanism for error reduction.

\textbf{Advantages:} Low variance; robust to overfitting; provides feature importance estimates;
parallelisable.

\textbf{Limitations:} Less interpretable than a single tree; higher memory and inference cost;
feature importances biased toward high-cardinality features.

% ─────────────────────────────────────────────────────────────────────────────
\section{Experimental Setup}
% ─────────────────────────────────────────────────────────────────────────────

\begin{tabular}{ll}
\toprule
\textbf{Component}     & \textbf{Detail} \\
\midrule
Python Version         & 3.12 \\
Scikit-learn Version   & 1.5.x \\
NumPy Version          & 1.26.x \\
Pandas Version         & 2.2.x \\
Matplotlib / Seaborn   & 3.9.x / 0.13.x \\
IDE                    & Jupyter Notebook / VS Code \\
Operating System       & Windows \\
Hardware               & CPU: Intel Core i5 (8-core), RAM: 16 GB \\
\bottomrule
\end{tabular}

% ─────────────────────────────────────────────────────────────────────────────
\section{Hyperparameter Selection via 5-Fold Cross-Validation}
% ─────────────────────────────────────────────────────────────────────────────

Hyperparameters were selected by exhaustive grid search over the spaces below.
Each combination was evaluated using \texttt{StratifiedKFold(n\_splits=5, shuffle=True, random\_state=42)},
and the configuration maximising average CV accuracy was retained.

\subsection*{8.1 Decision Tree — Search Space}

\begin{center}
\begin{tabular}{ll}
\toprule
\textbf{Hyperparameter}  & \textbf{Candidates Explored} \\
\midrule
\texttt{criterion}         & \texttt{gini}, \texttt{entropy} \\
\texttt{max\_depth}        & 3, 5, 7, \texttt{None} \\
\texttt{min\_samples\_split} & 2, 5, 10 \\
\texttt{min\_samples\_leaf}  & 1, 2, 4 \\
\bottomrule
\end{tabular}
\end{center}

\subsection*{8.2 Decision Tree Hyperparameter Evaluation (Top 8 by CV Accuracy)}

\begin{center}
\begin{tabular}{ccccccc}
\toprule
\textbf{Criterion} & \textbf{Max Depth} & \textbf{Avg CV Acc (\%)} & \textbf{Avg CV F1} \\
\midrule
gini      & 5    & 93.85 & 0.9506 \\
gini      & 5    & 93.85 & 0.9506 \\
gini      & 5    & 93.85 & 0.9506 \\
gini      & 7    & 93.85 & 0.9506 \\
gini      & None & 93.85 & 0.9506 \\
gini      & 7    & 93.41 & 0.9470 \\
gini      & 7    & 93.41 & 0.9470 \\
gini      & None & 93.41 & 0.9470 \\
\bottomrule
\end{tabular}
\end{center}

\textbf{Selected DT Configuration:}
\texttt{criterion=gini, max\_depth=5, min\_samples\_split=2, min\_samples\_leaf=4}

\subsection*{8.3 Random Forest — Search Space}

\begin{center}
\begin{tabular}{ll}
\toprule
\textbf{Hyperparameter}  & \textbf{Candidates Explored} \\
\midrule
\texttt{n\_estimators}   & 50, 100, 200 \\
\texttt{max\_depth}      & 5, 10, \texttt{None} \\
\texttt{max\_features}   & \texttt{sqrt}, \texttt{log2}, \texttt{None} \\
\texttt{bootstrap}       & \texttt{True}, \texttt{False} \\
\bottomrule
\end{tabular}
\end{center}

\subsection*{8.4 Random Forest Hyperparameter Evaluation (Top 8 by CV Accuracy)}

\begin{center}
\begin{tabular}{cccccc}
\toprule
\textbf{n\_est} & \textbf{Max Depth} & \textbf{Max Features}
  & \textbf{Avg CV Acc (\%)} & \textbf{Avg CV F1} \\
\midrule
100 & 10   & log2 & 96.48 & 0.9718 \\
100 & None & log2 & 96.48 & 0.9718 \\
50  & 10   & log2 & 96.26 & 0.9702 \\
50  & None & log2 & 96.26 & 0.9702 \\
100 & 10   & sqrt & 96.26 & 0.9699 \\
100 & None & sqrt & 96.26 & 0.9699 \\
200 & 10   & sqrt & 96.26 & 0.9700 \\
200 & 10   & log2 & 96.26 & 0.9703 \\
\bottomrule
\end{tabular}
\end{center}

\textbf{Selected RF Configuration:}
\texttt{n\_estimators=100, max\_depth=10, max\_features='log2', bootstrap=True}

% ─────────────────────────────────────────────────────────────────────────────
\section{5-Fold Cross-Validation Performance}
% ─────────────────────────────────────────────────────────────────────────────

\begin{center}
\begin{tabular}{lcccccc}
\toprule
\textbf{Model} & \textbf{Fold 1} & \textbf{Fold 2} & \textbf{Fold 3}
               & \textbf{Fold 4} & \textbf{Fold 5} & \textbf{Average} \\
               & (\%) & (\%) & (\%) & (\%) & (\%) & (\%) \\
\midrule
Decision Tree  & 94.51 & 94.51 & 91.21 & 92.31 & 96.70 & 93.85 \\
Random Forest  & 96.70 & 97.80 & 93.41 & 95.60 & 98.90 & 96.48 \\
\bottomrule
\end{tabular}
\end{center}

\begin{figure}[H]
\centering
\includegraphics[width=0.82\linewidth]{fig7_cv_folds}
\caption{Per-fold cross-validation accuracy for the best Decision Tree and Random Forest
  configurations. Dashed lines indicate each model's mean CV accuracy.}
\end{figure}

\textbf{Inference:}
Random Forest consistently outperforms the single Decision Tree across all five folds, with a mean
accuracy of 96.48\% versus 93.85\% for the Decision Tree. The Random Forest also shows lower
fold-to-fold variability (range: 5.49 pp vs. 5.49 pp), confirming that ensemble aggregation reduces
prediction variance. The largest gap occurs at Fold 5 (98.90\% vs. 96.70\%), where the forest's
diversity across 100 trees provides the greatest advantage.

% ─────────────────────────────────────────────────────────────────────────────
\section{Results}
% ─────────────────────────────────────────────────────────────────────────────

\subsection*{10.1 Test-Set Evaluation Metrics (80--20 Stratified Split)}

\begin{center}
\begin{tabular}{lcc}
\toprule
\textbf{Metric}   & \textbf{Decision Tree} & \textbf{Random Forest} \\
\midrule
Accuracy (\%)     & 90.35 & 95.61 \\
Precision (\%)    & 94.20 & 95.89 \\
Recall (\%)       & 90.28 & 97.22 \\
F1-Score (\%)     & 92.20 & 96.55 \\
AUC-ROC           & 0.9358 & 0.9924 \\
\bottomrule
\end{tabular}
\end{center}

\vspace{4pt}
Random Forest exceeds the Decision Tree on every metric. The most clinically significant gain is in
\textbf{Recall} (97.22\% vs. 90.28\%)---fewer malignant tumours are missed---and
\textbf{AUC-ROC} (0.9924 vs. 0.9358), indicating far better probability calibration across all
classification thresholds.

% ─────────────────────────────────────────────────────────────────────────────
\section{Result Visualizations}
% ─────────────────────────────────────────────────────────────────────────────

\subsection*{11.1 Confusion Matrix — Decision Tree}

\begin{figure}[H]
\centering
\includegraphics[width=0.52\linewidth]{fig3_dt_cm}
\caption{Confusion matrix for the Decision Tree classifier on the held-out test set (114 samples).}
\end{figure}

\textbf{Inference:}
The Decision Tree correctly classifies 38 of 42 malignant samples (TN) and 65 of 72 benign
samples (TP). It produces 7 false negatives (malignant samples predicted as benign)---the most
critical error type in oncology---and 4 false positives. The false negative rate of 16.7\% for
the malignant class indicates residual variance in the single-tree boundary.

\subsection*{11.2 Confusion Matrix — Random Forest}

\begin{figure}[H]
\centering
\includegraphics[width=0.52\linewidth]{fig4_rf_cm}
\caption{Confusion matrix for the Random Forest classifier on the held-out test set (114 samples).}
\end{figure}

\textbf{Inference:}
The Random Forest reduces false negatives from 7 to 2 compared to the Decision Tree, achieving a
malignant-class recall of 95.24\% (39/41). False positives drop from 4 to 3. The improvement in
true negatives (malignant correctly identified: 39 vs.\ 38) and dramatic reduction in missed
malignancies confirms that ensemble aggregation directly addresses the high-variance weakness of
individual trees on minority-class boundary regions.

\subsection*{11.3 ROC Curve Comparison}

\begin{figure}[H]
\centering
\includegraphics[width=0.65\linewidth]{fig5_roc}
\caption{ROC curves for the Decision Tree and Random Forest classifiers.
  Area under the curve (AUC) values are shown in the legend.}
\end{figure}

\textbf{Inference:}
The Random Forest achieves an AUC of 0.9924 compared to 0.9358 for the Decision Tree, confirming
near-perfect discrimination between malignant and benign classes across all probability thresholds.
The Decision Tree's curve shows a more step-like profile, reflecting its discrete class-probability
estimates (leaves assign the majority class probability uniformly). The Random Forest's smooth,
tightly hugging curve results from averaging probability estimates across 100 trees, producing
well-calibrated soft predictions throughout the operating range.

\subsection*{11.4 Feature Importances — Random Forest}

\begin{figure}[H]
\centering
\includegraphics[width=0.88\linewidth]{fig6_feat_imp}
\caption{Random Forest feature importances (mean decrease in impurity) for the top 15 features.
  Darker bars indicate features above the 0.08 importance threshold.}
\end{figure}

\textbf{Inference:}
The three most discriminative features are \textit{worst concave points} (0.1306),
\textit{worst area} (0.1186), and \textit{worst radius} (0.1071), all belonging to the
``worst'' measurement group which captures the most severe cell nucleus morphology in each image.
Features related to nuclear concavity and size dominate the model's decision space, consistent with
medical knowledge that malignant cells exhibit irregular, enlarged nuclei. The cumulative importance
of the top 10 features accounts for approximately 84\% of the total predictive weight, suggesting
the 30-feature space contains significant redundancy.

% ─────────────────────────────────────────────────────────────────────────────
\section{Discussion}
% ─────────────────────────────────────────────────────────────────────────────

\textbf{Tree depth and overfitting:}
In this experiment, constraining \texttt{max\_depth=5} alongside \texttt{min\_samples\_leaf=4}
achieved the same CV accuracy as an unconstrained tree (93.85\%), confirming that beyond a depth
threshold the model memorises noise rather than learning signal. The regularised tree generalises
identically while being far more compact and interpretable.

\textbf{Most impactful hyperparameter:}
For the Decision Tree, \texttt{min\_samples\_leaf} had the greatest isolated impact: increasing
it from 1 to 4 consistently lifted CV accuracy by 1--2 pp across all depth and criterion
combinations. For the Random Forest, \texttt{max\_features='log2'} (vs.\ \texttt{sqrt} or
\texttt{None}) produced the strongest gains, as it enforces tighter feature de-correlation between
trees.

\textbf{Why ensemble learning improved performance:}
Random Forest's improvement over the single tree is attributable to variance reduction through
two orthogonal mechanisms: (1) bootstrap aggregation distributes the learning signal across
different training subsets; (2) random feature subsets at each split de-correlate individual tree
errors. The majority vote then cancels these errors, cutting test error from 9.65\% to 4.39\%.

\textbf{Limitations:}
The Random Forest is less interpretable than a single decision tree and incurs higher inference
cost. Feature importance scores derived from mean decrease in impurity are biased toward
high-cardinality features and correlated variables (several here, e.g.\ radius/perimeter/area
triplets). Permutation importance would give a more reliable ranking. The dataset is also
relatively small (569 samples); performance on larger, noisier FNA cohorts should be validated.

Decision Tree and Random Forest classifiers were implemented on the Wisconsin Diagnostic Breast
Cancer dataset and tuned via exhaustive 5-fold stratified cross-validation over their respective
hyperparameter grids. The optimal Decision Tree (\texttt{criterion=gini, max\_depth=5,
min\_samples\_leaf=4}) achieved a mean CV accuracy of 93.85\% and a test-set accuracy of 90.35\%.
The optimal Random Forest (\texttt{n\_estimators=100, max\_depth=10, max\_features=log2,
bootstrap=True}) achieved a mean CV accuracy of 96.48\% and a test-set accuracy of 95.61\%,
outperforming the single tree across all five evaluation metrics including a notably higher
AUC-ROC (0.9924 vs.\ 0.9358) and recall for the malignant class (97.22\% vs.\ 90.28\%).

The results confirm that Random Forest's bootstrap aggregation and random feature sub-spacing
effectively reduce the high variance inherent in single decision trees, yielding more stable
and clinically reliable breast cancer classification.

% ─────────────────────────────────────────────────────────────────────────────
\section{References}
% ─────────────────────────────────────────────────────────────────────────────

\begin{enumerate}[label={[\arabic*]},leftmargin=*]

\item F. Pedregosa \textit{et al.}, ``Scikit-learn: Machine Learning in Python,''
  \textit{Journal of Machine Learning Research}, vol.\ 12, pp.\ 2825--2830, 2011.
  [Online]. Available: \url{https://scikit-learn.org}

\item L. Breiman, ``Random Forests,'' \textit{Machine Learning}, vol.\ 45, no.\ 1,
  pp.\ 5--32, Oct. 2001.
  doi: \href{https://doi.org/10.1023/A:1010933404324}{10.1023/A:1010933404324}

\item L. Breiman, J. Friedman, R. Olshen, and C. Stone,
  \textit{Classification and Regression Trees}. Belmont, CA: Wadsworth, 1984.

\item D. Dua and C. Graff, ``UCI Machine Learning Repository,''
  University of California, Irvine, School of Information and Computer Sciences, 2019.
  [Online]. Available: \url{https://archive.ics.uci.edu/ml/datasets/Breast+Cancer+Wisconsin+(Diagnostic)}

\item W. H. Wolberg, W. N. Street, and O. L. Mangasarian,
  ``Breast cancer Wisconsin (diagnostic) data set,''
  UCI Machine Learning Repository, 1995.

\item Scikit-learn Documentation, ``Decision Trees,'' 2024.
  [Online]. Available: \url{https://scikit-learn.org/stable/modules/tree.html}

\item Scikit-learn Documentation, ``Ensemble Methods -- Forests of Randomized Trees,'' 2024.
  [Online]. Available:
  \url{https://scikit-learn.org/stable/modules/ensemble.html#forests-of-randomized-trees}

\end{enumerate}

\end{document}